% s-331-p331-m002; original PDF 331, 332, 333, 334, 335
\Aphorism{s-331-p331-m002}{381}
Wenn wir einen \Sperr{Conditionalsatz} aussprechen, so fällen wir eigentlich das Urtheil, dass einzelne Theile der vorliegenden Erfahrungswelt herausgegriffen und in irgend einem Sinne variirt werden könnten, ohne dass andere Theile mit variirt würden: d. h. wir fällen ein unsinniges Urtheil. Denn in Strenge wird durch jedes Einzelphänomen das mit ihm durchweg verknüpfte kosmische Gesammtphänomen mitgesetzt, vorausgesetzt, reproducirt; kein Sperling fällt vom Dache, ohne das Universum in mechanische Mitleidenschaft zu ziehen, und zwar ebensowohl in aufsteigender wie absteigender Zeitrichtung. Die Welt ist allenfalls als Ganzes, ihrer Existenz oder Nichtexistenz nach, noch beliebig und frei; hat sich diese Alternative einmal entschieden, so giebt es innerhalb der Welt kein Wenn und Aber mehr, kein „hätte können“ und „hätte nicht sollen“, nichts wegzudenken und nichts hinzuzudenken. Die Unzulässigkeit conditionaler Urtheile wird nur dadurch zu einer \Sperr{bedingten} \Sperr{Zulässigkeit}, dass wir in die Betrachtung ein \Sperr{metrisches} Element und damit die Erlaubniss zu Annäherungen hineinbringen. Wir fassen jene Variationen (die entweder alle Null oder alle von Null verschieden sind) als extensive Grössen und setzen diejenigen von ihnen, die einen gewissen Betrag nicht übersteigen, practisch der Null gleich; wir urtheilen, dass es Deformationen der bestehenden Welt von derartiger Zusammensetzung geben könne, dass ein Theil über, ein anderer Theil unter einer gewissen \Sperr{Genauigkeitsgrenze} (etwa der menschlichen Bewusstseinsschwelle oder, gröber, der mindesten socialen Tragweite) vor sich ginge. Wir vernachlässigen also, mathematisch zu reden, die Glieder höherer Ordnung; und solange wir über den asymptotischen Charakter der hieraus fliessenden Ergebnisse im Klaren sind, wer will uns dieses Grundvorrecht alles organischen Bewusstseins streitig machen? Was sollte aus uns armen Erkennenden werden, wenn wir durchweg mit unbegrenzt vielen Decimalstellen rechnen müssten, ohne je abbrechen zu dürfen, wenn jeder Sperling, der auf irgendwelchem Siriusplaneten vom Dache fällt, nicht nur nominell in unserem Bewusstsein vertreten wäre, sondern seinen bestimmten Realwerth oberhalb der Reizschwelle hätte, wenn wir statt des geradlinigen Lichtstrahls Billionen Ätherschwingungen von fortwährend wechselnder Amplitude und Polarisationsebene wahrnehmen müssten? Man denke einmal an den Grenzfall, dass unsere Reizschwelle ganz und gar Null, unsere Reizhöhe ganz und gar unendlich wäre; so wie wir jetzt glücklicherweise organisirt sind, können wir nicht sehr tief in diese Möglichkeit hineinblicken, ohne Gleichgewicht, Athem und Besinnung zu verlieren. Ein Bewusstsein, das auf alles unterschiedslos antworten, das dem verzwickten Hin und Her zwischen erster Ursache, erster Wirkung, dadurch modificirter zweiter Ursache und entsprechend modificirter zweiter Wirkung und so fort bis ins Unendliche rastlos folgen muss, ein Bewusstsein, das, von aussen inducirt, „vorstellt“ und durch diese Vorstellung sowohl sich selbst, als rückwärts die Aussendinge und durch deren Vermittlung sich selbst noch einmal inducirt, ein Bewusstsein, dem sich alles zu Molecularschwärmen, Wirbelfeldern und Oscillationen auflöst, ein Ahasver-Bewusstsein, das nicht leben und nicht sterben kann — ist eine ärgere Marter denkbar als mit einem Bewusstsein dieser Art behaftet zu sein? Unter solchen Bedingungen werden die Elementarbegriffe wie Zeit, Leben, Tod, für uns illusorisch; denn was wir hierunter verstehen, die Verschiebung der Perspectiven, das unablässige Schwanken der Bewusstseinswelt nach Umfang, Tiefe und Gliederung, die Verengerung, Erweiterung, schliessliche Coincidenz der beiden Reizgrenzen — das alles hört auf, wenn diese Reizgrenzen ihre extremen Werthe angenommen haben. Man wird, mit der gehörigen Vorsicht in derlei Grenzbetrachtungen, einer solchen Bewusstseinswelt selbst die \Sperr{Zeit} als Grundform abzusprechen haben, gesetzt selbst, dass das Urbild, dessen vollständiges Abbild sie ist, objectiv zeitlich verliefe. Der Intellect hat nur an seinem eigenen Heller- und Trüber-, Stärker- und Schwächer-, Weiter- und Engerwerden die Vergleichsscala, längs deren er die Wahrnehmungen der Aussenwelt zeitlich auseinanderlegt; fehlt ihm diese, muss er als absoluter Spiegel die Umwelt vollständig und ohne Abzug, also mit samt der ganzen Vergangenheit und Zukunft, wiedergeben, so wird er — vermuthlich! — eine Art Nunc stans empfinden, d. h. nach unserem Maasse gemessen: überhaupt nicht empfinden. Kurzum: unser Bewusstsein ist, wesentlich und unvermeidlich, ein vernachlässigendes, vergröberndes, annäherndes, aufsaugendes und ausscheidendes Bewusstsein, denselben Gesetzen der Ernährung, Entwicklung und Anpassung unterworfen wie jede andere organische Function; nur insofern es endlich und begrenzt ist, insofern es sich das Recht nimmt, irgendwo abzubrechen, abzurunden und Halt zu machen, ist es überhaupt als Bewusstsein möglich. Der \Foreign{regressus in infinitum} wäre, wie Kant ausgeführt hat, eine Art Paralysis und Selbstmord der Erkenntniss, ein nicht zu beseitigender Widerspruch zwischen der unendlichen Reihe und ihrem abgeschlossenen Resultat, aber ein Widerspruch, der aus der blossen Speculation in die unmittelbare Anschauung verlegt wäre. — Somit haben, um zum Ausgang zurückzukehren, auch hypothetische Urtheile und die damit zusammenhängenden Begriffe wie Möglichkeit, freie Wahl, Verantwortlichkeit, ihren guten Sinn, nämlich als approximative Schätzungen von der Form: ich kann mir eine Welt denken, deren durchgängige Abweichung von der wirklichen stellenweise als merklich, stellenweise als unmerklich anzusehen ist; da wo zum ersten Male merkliche Änderungen auftreten, hat, nach unseren Begriffen, Freiheit und Zurechnungsfähigkeit ihren Sitz. Wir imputiren dem Verbrecher seine That, wir sagen, er hätte sie unterlassen können (und sollen), d. h. wir schätzen, dass keine Ummodelung des ganzen Weltverlaufs bis zum Zeitpunkte des Delicts erforderlich war, um zu diesem Zeitpunkte in eine Unterlassung des Delicts auszumünden. Das ist in Strenge unrichtig; es war eine solche Ummodelung erforderlich, aber sie kann vielleicht als verschwindend klein gelten. Vielleicht! — jene Schätzungen mögen im besonderen Falle falsch, ungenügend, oberflächlich sein; wir werden oft genug den \Sperr{Zeitpunkt}, die Trennungsstelle zwischen unmerklichen und merklichen Variationen verfehlen und den Verbrecher haftbar machen, wo vielleicht seine Eltern oder Vorgenerationen die Schuldigen waren. Aber die Grundvoraussetzung, dass überhaupt metrisch geschätzt, das Eine berücksichtigt, das Andere vernachlässigt werden dürfe, ist unanfechtbar; gehört sie nicht zu den Rechtmässigkeiten der Erkenntniss, so doch zu den Nothwendigkeiten des Lebens, und in diesem Falle heisst es \Foreign{primum vivere, deinde philosophari}!
\par

